\chapter{The Richards equation, roots and the Daisy ensemble}
\markboth{Richards equation, roots and Daisy}{Richards equation, roots and Daisy}
\textbf{After this chapter you can} write the one-dimensional Richards equation with a root-uptake sink, describe how Daisy closes the water balance, and read the design of the 10,190-run ensemble.

\section{Flow in an unsaturated soil}
Darcy's law for vertical flow, with $z$ positive upward and $h$ the suction (so that the pressure head is $-h$), gives the flux
\begin{equation}\label{eq:darcy}
q = -K(h)\left(\frac{\partial(-h)}{\partial z}+1\right)=K(h)\left(\frac{\partial h}{\partial z}-1\right).
\end{equation}
Conservation of mass for the water content $\theta$ with a sink $S_r$ for root uptake gives
\begin{equation}\label{eq:richards}
\frac{\partial\theta}{\partial t}=-\frac{\partial q}{\partial z}-S_r=\frac{\partial}{\partial z}\left[K(h)\left(1-\frac{\partial h}{\partial z}\right)\right]-S_r .
\end{equation}
Equation~(\ref{eq:richards}) is the Richards equation. Together with the retention curve $\theta(h)$ of Equation~(\ref{eq:vg}) and the conductivity of Equation~(\ref{eq:vgm}) it is a closed, non-linear equation for $h(z,t)$. It solves the \emph{equilibrium} form: the water content follows the retention curve at once. When wetting and drying are fast, the water content lags behind, and the criterion of Section~\ref{sec:lag} says when this matters.

\section{Uptake by roots: the conductance of the soil around a root}
Roots take water from the soil around them. For a cylindrical root of radius $r_0$ in a soil cylinder of radius $r_1$, steady radial flow gives an uptake per unit root length
\begin{equation}\label{eq:radial}
u\simeq\frac{2\pi}{\ln(r_1/r_0)}\int_{h_{\mathrm{root}}}^{h}K(h')\,dh'\approx\frac{2\pi}{\ln(r_1/r_0)}\,K(h)\,\Delta h ,
\end{equation}
where $\Delta h$ is the head drop from the bulk soil to the root surface (Gardner, 1960). Multiplying by the root length density $L_v$ gives the sink per unit volume. Two things follow. First, the uptake is proportional to $K(h)$, so as the soil dries and $K$ falls by orders of magnitude (Chapter~3), uptake falls with it. Second, $\Delta h$ is bounded by the plant, since the root water potential cannot fall below the wilting point. So there is a suction beyond which the soil delivers too little water for the crop to dry it further in a season. This is the origin of the reach (Chapter~\ref{ch:reach}). Daisy computes uptake from a soil--root potential difference over a conductance of this type (Abrahamsen and Hansen, 2000).

\begin{example}{the collapse of uptake}
Take $\Delta h$ and the geometry as fixed and compare the uptake at two suctions in the silty clay loam S10. From Chapter~3, $\Kr(300\text{ cm})=1.5\times10^{-3}$ and $\Kr(3000\text{ cm})\approx 5.5\times10^{-6}$ (from the asymptote of Chapter~\ref{ch:reach}). With the same $\Delta h$ and root length density the uptake falls by a factor of about 280 when the soil dries from 300 to 3000 cm. In the sand the same drying takes $\Kr$ from $2.7\times10^{-11}$ to about $2\times10^{-17}$: the uptake is already negligible at 300 cm.
\end{example}

\section{What Daisy is}
Daisy is a soil, plant and atmosphere model developed at the University of Copenhagen. The version used here is 7.1.14. It solves Equation~(\ref{eq:richards}) in a layered one-dimensional column with the functions above. It couples the soil to a crop (light interception, photosynthesis, roots, water stress), to a surface residue layer that intercepts rain and suppresses soil evaporation, to an energy-based potential evapotranspiration, and to nitrogen turnover. For every run it writes daily logs of water content, suction, the water fluxes, the crop and the nitrogen pools.

\fig{sh_dzmap}{0.95}{The Daisy process map as shown in the SWAFL workbench. Clicking a box opens the description of the process and its figures.}

\section{The ensemble}
The ensemble has 17 experiments, E0 to E16, ten soils (Table~\ref{tab:soils}), 20 years of weather from Taastrup, Denmark (1980--1999) and a cereal rotation. Of 10,190 runs 10,041 are valid and 149 are solver failures, which are excluded. The irrigation scheduler of Daisy applies a dose whenever the suction at a sensor depth exceeds a trigger. Every run is archived: the one-row-per-run summary, every log year by year, and the summariser are in the repository, and the complete raw archive is on Zenodo (doi 10.5281/zenodo.23043626).

\begin{table}[H]\centering\small\caption{Experiments used in this book.}\begin{tabular}{llp{6.4cm}r}\toprule Exp.&Name&Factors varied&Valid runs\\\midrule
E0&Baselines&soil (the rainfed twin of every run)&10\\
E2&TriggerPressure&soil, trigger from 50 to 10,000 cm; sensor at 20 cm&197\\
E3&SensorDepth&soil, sensor depth 5 to 100 cm at a 300 cm trigger&200\\
E4&Methods&soil, method (surface, overhead, subsurface)&96\\
E9&Initial suction lattice&soil, initial suction (10 values), rainfall scale (10 values)&987\\
E10&Amount--timing--soil&soil, dose (1 to 30 mm), irrigation window&954\\
E15&Method--depth--amount&coarse loam; method, rate, dose, rooting depth&1000\\
E16&Factorial&soil; six factors at two levels&740\\\bottomrule\end{tabular}\end{table}

\begin{daisybox}{reading an annual flux}
For each run and each year the summariser adds the daily fluxes of transpiration, soil evaporation, other evaporation, drainage at 1 m and irrigation. A twenty-year mean per run is what the equations of the next chapters use as $T$, $E_s$, $E_o$, $D$ and $I$. The other evaporation is $E_o=ET-T-E_s$, and in the data it closes the evapotranspiration to within 1.1 mm.
\end{daisybox}

\section{Sensor lag: when the equilibrium assumption matters}\label{sec:lag}
Under fast wetting and drying the water content lags the equilibrium of the retention curve, $\partial\theta/\partial t=[\theta_{\mathrm{eq}}(h)-\theta]/\tau$. A sensor then reads a suction that is not the one of Equation~(\ref{eq:vg}). For a schedule with interval $t_i$ between applications, the dimensionless lag $\Pi_\tau=\tau/t_i$ says when this matters. For daily or multi-day scheduling ($t_i\ge1$ d) and $\tau$ of minutes to hours, $\Pi_\tau\ll1$ and the equilibrium reach applies. For pulses of minutes to hours, as in pulsed drip, $\Pi_\tau\sim1$ and the sensor reading becomes rate dependent. Daisy solves the equilibrium equation, so this criterion is stated here for later testing; it is not evaluated in the studies.

\begin{exercise}A schedule irrigates every 2 days and the sensor lag is $\tau=3$ h. Compute $\Pi_\tau$ and say whether the equilibrium reach applies. Repeat for pulses every 1 h.\end{exercise}
\begin{exercise}From Equation~(\ref{eq:radial}), show that if $K$ falls by a factor of $10^3$ and $\Delta h$ is fixed, the uptake falls by the same factor.\end{exercise}
\begin{exercise}Starting from Equation~(\ref{eq:darcy}), derive Equation~(\ref{eq:richards}) by applying conservation of mass to a thin layer.\end{exercise}
\begin{exercise}Why does Daisy's solution of the equilibrium Richards equation not need a sensor-lag parameter?\end{exercise}

\chapter{The reach of a suction trigger}\label{ch:reach}
\markboth{The reach of a suction trigger}{The reach of a suction trigger}
\textbf{After this chapter you can} derive the dry-end asymptote of the conductivity, derive the closed form for the reach, compute the reach for any soil with known $\alpha$ and $n$, and explain why the constant $\kappa$ is the same for all soils.

\section{The question}
Sensor-scheduled irrigation applies water when the suction at a sensor passes a trigger $h_t$. Triggers of 300 to 600 cm come from tables, the same for every soil. But a soil may never reach such a suction. If the crop cannot dry the soil beyond some suction $\hc$, then a trigger $h_t>\hc$ is never reached, the crop is not irrigated, and the field is in effect rainfed. A trigger $h_t<\hc$ can fire. The quantity $\hc$ is the \emph{reach} of the trigger.

\section{The dry-end asymptote}
Far from saturation, $\alpha h\gg1$. Then $\Se^{1/m}=[1+(\alpha h)^n]^{-1}\approx(\alpha h)^{-n}$ and, since $m n=n-1$, $\Se\approx(\alpha h)^{-(n-1)}$. In Equation~(\ref{eq:vgm}), $1-(1-y)^m\approx m y$ for small $y=\Se^{1/m}$, so the bracket is $\approx m(\alpha h)^{-n}$ and
\begin{equation}\label{eq:asym}
\boxed{\;\Kr(h)\simeq m^{2}(\alpha h)^{-[\,2n+\ell(n-1)\,]},\qquad\alpha h\gg1.\;}
\end{equation}
The conductivity is a power law of suction, and the exponent $E=2n+\ell(n-1)$ grows with $n$.

\begin{example}{how good is the asymptote?}
For the silty clay loam S10 ($\alpha=0.010$, $n=1.23$, $m=0.187$, $E=2.575$): at $h=1000$ cm the exact value is $\Kr=8.64\times10^{-5}$ and Equation~(\ref{eq:asym}) gives $9.30\times10^{-5}$, an error of 8 per cent. At $h=10{,}000$ cm the exact value is $2.46\times10^{-7}$ and the asymptote $2.48\times10^{-7}$, an error under 1 per cent. At 100 cm, where $\alpha h=1$, the asymptote $3.5\times10^{-2}$ is wrong by a factor of 2.5, because the assumption $\alpha h\gg1$ fails. The asymptote is therefore reliable at the suctions where triggers sit.
\end{example}

\section{One constant for all soils}
Suppose the crop stops drying the soil when the relative conductivity has fallen to a value $\kappa$. The justification is Chapter 4: the uptake collapses once $K$ is small enough. Setting $\Kr(\hc)=\kappa$ in Equation~(\ref{eq:asym}),
\begin{equation}\label{eq:reach}
\boxed{\;\hc=\frac1\alpha\left(\frac{m^{2}}{\kappa}\right)^{1/[2n+\ell(n-1)]}\;}
\end{equation}
The constant $\kappa$ describes the crop, the climate and the sensor. It does not depend on the soil, which enters only through $\alpha$, $n$ and $\ell$. With $\kappa=3.5\times10^{-10}$ and $\ell=0.5$:

\begin{table}[H]\centering\small\caption{Reach of a trigger from Equation~(\ref{eq:reach}), $\kappa=3.5\times10^{-10}$, $\ell=0.5$, computed from the rounded parameters of Table~\ref{tab:soils}; the study's values differ by under 1 per cent. The last column: does a 300 cm trigger fire?}\label{tab:reach}
\begin{tabular}{llrrrrl}\toprule ID&Texture&$n$&$E$&$\alpha\hc$&$\hc$ (cm)&300 cm\\\midrule
S03&Sand&2.680&6.200&28.8&199&never\\S04&Loamy sand&2.280&5.200&52.7&425&never\\S05&Sandy loam&1.890&4.225&121&1,615&fires\\S02&Coarse sand&1.600&3.500&287&5,359&fires\\S08&Sandy clay loam&1.480&3.200&446&7,559&fires\\S01&Coarse loam&1.353&2.882&751&7,814&fires\\S06&Loam&1.560&3.400&331&9,186&fires\\S07&Silt loam&1.410&3.025&591&29,528&fires\\S09&Clay loam&1.310&2.775&905&47,608&fires\\S10&Silty clay loam&1.230&2.575&1,278&127,830&fires\\\bottomrule\end{tabular}\end{table}
(The study reports values above 10,000 cm as bounded from below, because the trigger experiment stops at 10,000 cm.)

\begin{example}{the reach of a soil by hand}
A soil has $\alpha=0.05$ cm$^{-1}$ and $n=2.0$. Then $m=0.5$, $E=2(2.0)+0.5(1.0)=4.5$, $m^2/\kappa=0.25/(3.5\times10^{-10})=7.14\times10^{8}$, $\ln(7.14\times10^8)=20.39$, and $20.39/4.5=4.531$, so $(m^2/\kappa)^{1/E}=e^{4.531}=92.9$. The reach is $\hc=92.9/0.05=1{,}856$ cm. A trigger of 300 cm fires in this soil. In practice one keeps the trigger below about half the reach, here under 930 cm, because irrigation falls to zero as $h_t/\hc$ approaches 1 (figure at the end of this chapter).
\end{example}

\subsection{Properties of the reach}
\begin{enumerate}[leftmargin=*,itemsep=2pt]
\item $\alpha$ only rescales the reach: $\alpha\hc$ depends on $n$, $\ell$ and $\kappa$ alone.
\item \emph{Sensitivity to $n$.} Write $\ln(\alpha\hc)=\ln(m^2/\kappa)/E$. With $dm/dn=1/n^2$,
\begin{equation}\label{eq:sens}\frac{\partial\ln(\alpha\hc)}{\partial n}=\frac{2}{m\,n^2E}-\frac{(2+\ell)\ln(m^2/\kappa)}{E^2}.\end{equation}
The second term dominates, since $\ln(m^2/\kappa)$ is between 18 and 21. The reach therefore shortens steeply as $n$ grows: $\alpha\hc\approx1{,}280$ at $n=1.23$ and $\approx29$ at $n=2.68$.
\item A larger $\ell$ shortens the reach, because the conductivity collapses faster.
\item For $n\to1$ the exponent tends to 2 and $m\to0$, so the reach tends to $\alpha^{-1}(m^2/\kappa)^{1/2}$, which vanishes slowly. Very broad pore-size distributions remain reachable far into the dry range, as observed for the silty clay loam.
\item \emph{Robustness.} Because $\hc\propto\kappa^{-1/E}$, a tenfold change of $\kappa$ changes the reach by $10^{1/E}$: 1.45 in the sand ($E=6.2$) and 2.44 in the silty clay loam ($E=2.575$).
\end{enumerate}

\begin{keybox}A suction trigger has a reach $\hc=\alpha^{-1}(m^2/\kappa)^{1/[2n+\ell(n-1)]}$. A trigger beyond it can never fire. One constant, $\kappa\approx3.5\times10^{-10}$, serves all ten soils; left-out soils are predicted within a factor of 2.3.\end{keybox}

\section{Why a trigger has a reach, in words}
A crop dries the soil at the sensor only as fast as the soil delivers water to the roots. That delivery scales with the conductivity, which falls by many orders of magnitude as the soil dries. In a soil with a narrow pore-size distribution (large $n$, sands) the conductivity collapses within a narrow band of suction just past the air entry, so the crop stops drying the soil there, and whatever trigger is set beyond that band is never read. In a soil with a broad distribution (small $n$, silty clay loams) the conductivity fades slowly over many decades, and the crop can still dry the soil at very high suctions.

\section{Why a deeper sensor does not help}
Below the sensor depth of the experiments the root density is smaller, so the soil dries even less, and in sands the deeper soil is also wetted by drainage from above. A deeper sensor therefore reads lower suctions, not higher ones.

\section{The reach scale}
Define the reach-scaled trigger $x=h_t/\hc$. Plotting irrigation against $x$ instead of against $h_t$ collapses the curves of ten soils: the mean within-bin variance of the normalised irrigation, in ten quantile bins, falls from 54 per cent of the total variance (on $h_t$ alone) to 32 per cent (on $x$). It does not fall to zero, because the reach sets where irrigation ends, but how irrigation fades \emph{before} the reach depends on how often the soil reaches a given suction in a season, which depends on the water held between field capacity and that suction: a retention property, not a conductivity property.

\label{fig:reachscale}
\fig{p4_reach}{0.9}{Reach scale and closed form against Daisy (same figure as in Chapter~3, now read with Equation~(\ref{eq:reach})).}

\begin{exercise}Compute $\hc$ for $\alpha=0.03$ cm$^{-1}$, $n=1.7$ with $\kappa=3.5\times10^{-10}$ and $\ell=0.5$.\end{exercise}
\begin{exercise}Derive Equation~(\ref{eq:asym}) from Equation~(\ref{eq:vgm}), step by step, stating each approximation.\end{exercise}
\begin{exercise}Derive Equation~(\ref{eq:sens}). Evaluate it at $n=1.5$ with $\kappa=3.5\times10^{-10}$ and $\ell=0.5$ and say whether the reach rises or falls with $n$.\end{exercise}
\begin{exercise}For the soil of the worked example, by what factor does the reach change if $\kappa$ is $3.5\times10^{-9}$ instead of $3.5\times10^{-10}$?\end{exercise}
\begin{exercise}A farmer sets a 300 cm trigger in a loamy sand. Use Table~\ref{tab:reach} to explain why the field is in effect rainfed, and propose a trigger.\end{exercise}
